\documentclass[10pt,a4paper]{amsart}
\usepackage[margin=19mm]{geometry}
\usepackage{amsmath,amssymb,mathtools}
\usepackage[scr=rsfs,cal=euler]{mathalfa}
\usepackage{xfrac}
\usepackage[unicode,hidelinks,bookmarks=false]{hyperref}
\newcommand{\ie}{i.e.\ }
\newcommand{\cf}{cf.\ }
\newcommand{\resp}{respectively\ }
\newcommand{\Subsection}{\subsection}
\providecommand{\nobreakdash}{\nobreak\hbox{-}}
\setlength{\emergencystretch}{2em}
\multlinegap0pt
\usepackage{amssymb}
\usepackage{float}
\usepackage{mathtools}
\usepackage{bm}
\usepackage{xcolor}
\usepackage[shortlabels]{enumitem}
\newtheorem{lem}{Lemma}
\newcommand{\addmr}{\mathfrak{addmr}}
\DeclareMathOperator{\vimo}{vimo}
\title{A Lie algebra associated with adjoint multiple zeta values}
\author{Takumi Anzawa}
\date{Source: arXiv:2511.03177v3 (CC BY 4.0)}
\begin{document}
\maketitle

\noindent\emph{Source and licence.} A Lie algebra associated with adjoint multiple zeta values. Version arXiv:2511.03177v3 (CC BY 4.0), distributed by arXiv under the Creative Commons Attribution 4.0 International licence, http://creativecommons.org/licenses/by/4.0/, as recorded in the arXiv licence metadata for this exact version and shown on the abstract page https://arxiv.org/abs/2511.03177v3. Full licence notice: \texttt{licenses/arxiv-2511.03177-CC-BY-4.0.txt}.\par\smallskip

\noindent\textit{Editorial scope.} Counted unit: the article's lem lem (source0.tex lines 3285--3336). The article's complete original proof of it is delivered with it (source0.tex lines 3338--3486, one \texttt{proof} environment of 149 lines). No mathematics is written, repaired or re-derived: the statement and the proof are the article's own lines, in the article's own notation, and no retained line is substituted or re-flowed.  No further source-local context is needed: every cross-reference of every retained line resolves inside the two delivered spans.  Nothing was omitted from any retained span, so the package carries no omission record.  No retained line contains a citation, so the wrapper carries no bibliography and no bibliography entry.  No retained line contains a figure environment, an image-inclusion command or drawing code, so no image is delivered and the package declares no asset dependency.  Editorial formatting only, in addition: an AMS \texttt{amsart} preamble that restates the article's own declarations and macros used by the retained lines, namely the environments \texttt{lem} on separate counters, together with the standard packages those lines need.  The retained environments print in this document as Lemma 1 in order of appearance, whereas the article prints the same items with the numbering of its own full document.  The complete native source of the article as distributed in the arXiv e-print for this exact version is bundled unchanged as \texttt{sources/188483/source0.tex}.
\begin{lem}
Let $\Psi_1 \in \addmr_0 \cap F_2^{\operatorname{ad}(x_1)}$ and $\Psi_2 \in h^\vee$.
Then
\[
\vimo_{d_{\Psi_1}(\Psi_2)}(z,X*Y,z')
=
\mathcal{B}_{\Psi_1,\Psi_2}(z;X,Y;z')
+
\mathcal{X}_{\Psi_1,\Psi_2}(z;X,Y;z')
+
\mathcal{Y}_{\Psi_1,\Psi_2}(z;X,Y;z'),
\]
where
\[
\mathcal{B}_{\Psi_1,\Psi_2}(z;X,Y;z')
:=
A_{zX'}^{\mathrm{L}}+A_{zY'}^{\mathrm{L}}+A_{Xz'}^{\mathrm{R}}+A_{Yz'}^{\mathrm{R}}
-
A_{zX'}^{0}-A_{zY'}^{0},
\]
\[
\mathcal{X}_{\Psi_1,\Psi_2}(z;X,Y;z')
:=
A_{Xz'}^{\mathrm{L}}+A_{XX'}^{\mathrm{L}}+A_{XY'}^{\mathrm{L}}
+
A_{zX'}^{\mathrm{R}}+A_{XX'}^{\mathrm{R}}+A_{YX'}^{\mathrm{R}}
-
A_{Xz'}^{0}-A_{XX'}^{0}-A_{XY'}^{0},
\]
and
\[
\mathcal{Y}_{\Psi_1,\Psi_2}(z;X,Y;z')
:=
A_{Yz'}^{\mathrm{L}}+A_{YX'}^{\mathrm{L}}+A_{YY'}^{\mathrm{L}}
+
A_{zY'}^{\mathrm{R}}+A_{XY'}^{\mathrm{R}}+A_{YY'}^{\mathrm{R}}
-
A_{Yz'}^{0}-A_{YX'}^{0}-A_{YY'}^{0}.
\]
Here
\[
\vimo_{d_{\Psi_1}(\Psi_2)}(z,X*Y,z')
=
A_{zz'}+A_{zX'}+A_{zY'}+A_{Xz'}+A_{Yz'}+A_{XX'}+A_{XY'}+A_{YX'}+A_{YY'}
\]
is the decomposition according to the positions of the distinguished letters
appearing in Proposition~6.10, and for each $A_{\bullet}\neq A_{zz'}$,
\[
A_{\bullet}=A_{\bullet}^{\mathrm{L}}+A_{\bullet}^{\mathrm{R}}-A_{\bullet}^{0}
\]
denotes the decomposition obtained by applying \emph{(27)} to the $f_1$-factor.
\end{lem}

\begin{proof}
By Proposition~6.10, we have
\[
\vimo_{d_{\Psi_1}(\Psi_2)}(z,X*Y,z')
=
\sum_{V_L(v)V_M(v')V_R=(z)(X*Y)(z')}
f_1(v,V_M,v')\,
f_2(\delta_z(v),V_L,(v),(v'),V_R,\delta_{z'}(v')).
\]
We divide this sum according to the positions of $v$ and $v'$.
More precisely, there are nine cases:
\[
(z,z'),\quad
(z,X'),\quad
(z,Y'),\quad
(X,z'),\quad
(Y,z'),\quad
(X,X'),\quad
(X,Y'),\quad
(Y,X'),\quad
(Y,Y').
\]
This gives
\[
\vimo_{d_{\Psi_1}(\Psi_2)}(z,X*Y,z')
=
A_{zz'}+A_{zX'}+A_{zY'}+A_{Xz'}+A_{Yz'}+A_{XX'}+A_{XY'}+A_{YX'}+A_{YY'}.
\]

We write down two representative cases.

First, consider the case $(v,v')=(z,x')$ with $x'$ occurring in $X$.
Write $X=X_L(x')X_R$.
By Lemma~6.11,
\[
X*Y
=
\sum_{\substack{Y=Y_L(y)Y_R}}
c(x',y)\,(X_L*Y_L)\bigl((x')-(y)\bigr)(X_R*Y_R).
\]
Substituting this into the above formula and extracting the terms
corresponding to $(v,v')=(z,x')$, we obtain
\[
A_{zX'}
=
\sum_{\substack{X=X_L(x')X_R\\ Y=Y_L(y)Y_R}}
c(x',y)\,
f_1(z,X_L*Y_L,x')\,f_2(z,x',X_R*Y_R,z').
\]

Next, consider the case $(v,v')=(x,x')$ with both $x$ and $x'$ occurring in $X$.
Write $X=X_L(x)X_M(x')X_R$.
Applying Lemma~6.11 twice, we obtain
\[
X*Y
=
\sum_{\substack{Y=Y_L(y)Y_M(y')Y_R}}
c(x,y)c(x',y')\,
(X_L*Y_L)\bigl((x)-(y)\bigr)(X_M*Y_M)\bigl((x')-(y')\bigr)(X_R*Y_R).
\]
Hence the contribution of the case $(v,v')=(x,x')$ is
\[
A_{XX'}
=
\sum_{\substack{X=X_L(x)X_M(x')X_R\\ Y=Y_L(y)Y_M(y')Y_R}}
c(x,y)c(x',y')\,
f_1(x,X_M*Y_M,x')\,f_2(z,X_L*Y_L,x,x',X_R*Y_R,z').
\]

The remaining seven cases are obtained in exactly the same way.

Now, since $\Psi_1 \in \addmr_0$, Proposition~6.5 yields
\[
A_{zz'}=f_1(z,X*Y,z')f_2(z,z')=0.
\]
Since $\Psi_1 \in F_2^{\operatorname{ad}(x_1)}$, we may apply \emph{(27)} to the
$f_1$-factor in each of the remaining eight terms.
Thus, for every $A_{\bullet}\neq A_{zz'}$, we have
\[
A_{\bullet}=A_{\bullet}^{\mathrm{L}}+A_{\bullet}^{\mathrm{R}}-A_{\bullet}^{0},
\]
where
\[
A_{\bullet}^{\mathrm{L}}
\]
denotes the term obtained by replacing the factor
$f_1(a,W,b)$ by $f_1(a,W,0)$,
\[
A_{\bullet}^{\mathrm{R}}
\]
denotes the term obtained by replacing it by $f_1(0,W,b)$,
and
\[
A_{\bullet}^{0}
\]
denotes the term obtained by replacing it by $f_1(0,W,0)$.

Finally, we regroup these terms according to the position of the endpoint
which is replaced by $0$.

The terms
\[
A_{zX'}^{\mathrm{L}}+A_{zY'}^{\mathrm{L}}+A_{Xz'}^{\mathrm{R}}+A_{Yz'}^{\mathrm{R}}
-
A_{zX'}^{0}-A_{zY'}^{0}
\]
form the boundary contribution
\[
\mathcal{B}_{\Psi_1,\Psi_2}(z;X,Y;z').
\]

The terms
\[
A_{Xz'}^{\mathrm{L}}+A_{XX'}^{\mathrm{L}}+A_{XY'}^{\mathrm{L}}
+
A_{zX'}^{\mathrm{R}}+A_{XX'}^{\mathrm{R}}+A_{YX'}^{\mathrm{R}}
-
A_{Xz'}^{0}-A_{XX'}^{0}-A_{XY'}^{0}
\]
form the $X$-bulk contribution
\[
\mathcal{X}_{\Psi_1,\Psi_2}(z;X,Y;z').
\]

The terms
\[
A_{Yz'}^{\mathrm{L}}+A_{YX'}^{\mathrm{L}}+A_{YY'}^{\mathrm{L}}
+
A_{zY'}^{\mathrm{R}}+A_{XY'}^{\mathrm{R}}+A_{YY'}^{\mathrm{R}}
-
A_{Yz'}^{0}-A_{YX'}^{0}-A_{YY'}^{0}
\]
form the $Y$-bulk contribution
\[
\mathcal{Y}_{\Psi_1,\Psi_2}(z;X,Y;z').
\]

Therefore,
\[
\vimo_{d_{\Psi_1}(\Psi_2)}(z,X*Y,z')
=
\mathcal{B}_{\Psi_1,\Psi_2}(z;X,Y;z')
+
\mathcal{X}_{\Psi_1,\Psi_2}(z;X,Y;z')
+
\mathcal{Y}_{\Psi_1,\Psi_2}(z;X,Y;z'),
\]
as desired.
\end{proof}

\end{document}
